\section{相关工作}
\label{sec:related}

\paragraph{基于 CLIP 的零样本异常检测。}
WinCLIP~\cite{jeong2023winclip} 通过组合式状态提示与多尺度窗口特征建立语言引导的零样本分类与分割；AnomalyCLIP~\cite{zhou2024anomalyclip} 从辅助异常数据中学习物体无关的正常/异常文本提示，以提升未见类别上的泛化；VCP-CLIP~\cite{qu2024vcpclip} 将视觉上下文注入提示，AA-CLIP~\cite{ma2025aaclip} 则同时重塑文本锚点和视觉特征。这些方法主要回答“如何获得更适合异常检测的跨类别语义表示”。本文关注一个互补问题：当文本原型固定后，如何利用单张测试图的局部外观校准该原型，而不在目标类别上训练。

\paragraph{面向异常检测的频率方法。}
FE-CLIP~\cite{gong2025feclip} 用 DCT 和可学习适配器在 CLIP 视觉编码器的多层特征中注入频率信息。其他频率异常检测工作也常把高、低频分量作为互补特征分支。它们说明局部频谱结构对细粒度异常有价值，但频率通常直接进入表征学习或最终判别。本文不把小波分量作为新增特征分支，也不把频率图与异常图相加；小波只产生逐 patch 可靠性，用于控制哪些视觉 token 可以参与逐图原型构造。因此，本文的创新点不是“在 ZSAD 中首次使用频率”，而是明确频率在线语义参照构造中的角色。

\paragraph{测试时与参照式适配。}
WinCLIP+~\cite{jeong2023winclip} 从少量正常参考图中构建视觉关联分数；其他测试时方法会依赖伪标签、合成异常或外部记忆。本文既不访问目标类别正常库，也不在测试时优化参数，而是从当前图像中汇聚视觉原型。与纯语义自参照不同，小波可靠性为 patch 选择提供一条不由文本相似度直接产生的辅助依据。为避免过度声称，我们不宣称该范式绝对新颖，而是通过直接融合、纯语义证据与不同小波设计的受控比较来验证其必要性。
